\documentclass[11pt,a4paper]{article}
\usepackage[margin=1in]{geometry}
\usepackage{amsmath,amssymb,amsthm}
\usepackage{algorithm}
\usepackage{algpseudocode}
\usepackage{booktabs}
\usepackage{hyperref}

\theoremstyle{definition}
\newtheorem{definition}{\textbf{Definition}}[section]
\newtheorem{theorem}{\textbf{Theorem}}[section]
\newtheorem{lemma}[theorem]{\textbf{Lemma}}
\newtheorem{remark}{\textbf{Remark}}[section]

\title{\textbf{Redundancy-Reduction Objectives in Self-Supervised Learning}}
\author{\textbf{Team Scaibu} \\ \small Email: \texttt{jaydeep.9664920749@gmail.com} \\ \small \textbf{Architecture and Core Engine Team}}
\date{\textbf{October 2026}}

\begin{document}
\maketitle

\section{Definitions and Cognitive Mental Models}

\subsection{Formal Mathematical Definition}
\textbf{Definition (Barlow Twins Cross-Correlation and Redundancy Reduction):}
Let $\mathbf{Z}_A, \mathbf{Z}_B \in \mathbb{R}^{B \times D}$ denote the batch representations of two augmented views generated by an encoder-projector network for a mini-batch of size $B$. Each feature column is zero-centered and normalized to unit variance along the batch dimension:
{\boldmath
\begin{equation}
\tilde{Z}_{A, bi} = \frac{Z_{A, bi} - \mu_{A, i}}{\sigma_{A, i} + \epsilon}, \quad \tilde{Z}_{B, bi} = \frac{Z_{B, bi} - \mu_{B, i}}{\sigma_{B, i} + \epsilon}
\end{equation}
}
The empirical \textbf{Cross-Correlation Matrix} $\mathbf{C} \in \mathbb{R}^{D \times D}$ between the two representations is defined as:
{\boldmath
\begin{equation}
C_{ij} = \frac{1}{B} \sum_{b=1}^B \tilde{Z}_{A, bi} \, \tilde{Z}_{B, bj}
\end{equation}
}
The \textbf{Barlow Twins Objective} forces $\mathbf{C}$ toward the identity matrix $\mathbf{I}_D$:
{\boldmath
\begin{equation}
\mathcal{L}_{\text{Barlow}} = \underbrace{\sum_{i=1}^D (1 - C_{ii})^2}_{\text{Invariance Term}} + \lambda \underbrace{\sum_{i=1}^D \sum_{j \ne i}^D C_{ij}^2}_{\text{Redundancy Reduction Term}}
\end{equation}
}
where $\lambda > 0$ balances the penalty between feature invariance and off-diagonal cross-correlation decorrelation.

Under the related \textbf{VICReg} formulation, collapse is prevented via three explicit decoupled criteria:
{\boldmath
\begin{equation}
\mathcal{L}_{\text{VICReg}} = \lambda \, s(\mathbf{Z}_A, \mathbf{Z}_B) + \mu \left( v(\mathbf{Z}_A) + v(\mathbf{Z}_B) \right) + \nu \left( c(\mathbf{Z}_A) + c(\mathbf{Z}_B) \right)
\end{equation}
}
where $s$ enforces mean squared error invariance, $v$ penalizes feature variance falling below a threshold $\gamma$, and $c$ penalizes off-diagonal covariance entries.

\subsection{Expanded Non-Technical Definition (Intuition for Anyone)}
\textbf{Intuitive Conceptual Definition:}
\begin{enumerate}
    \item \textbf{Horace Barlow's 1961 Brain Hypothesis}: Neuroscientist Horace Barlow discovered that when biological eyes process vision, individual brain neurons try to be as independent as possible. If two neurons are sending the exact same message, one of them is completely redundant and wasting precious metabolic energy!
    \item \textbf{No Negatives, No Clones, Just Teamwork}: Barlow Twins doesn't need thousands of negative images, nor does it need weird stop-gradient tricks. It treats the 2,048 neurons in an AI embedding vector like a team of specialized detectives.
    \item \textbf{Rule 1 (Agree on the Clues)}: If Detective 5 in view 1 spots a red car, Detective 5 in view 2 must also spot the red car (diagonal cross-correlation equals 1).
    \item \textbf{Rule 2 (Don't Copy Your Colleague)}: Detective 5 must not report what Detective 6 is already reporting. If Detective 6 covers shapes, Detective 5 must cover colors. Every off-diagonal correlation must equal 0.
    \item \textbf{Collapse is Physically Impossible}: Because every single neuron is forced to tell a completely different, independent story about the photo, all 2,048 neurons cannot collapse into a boring constant.
\end{enumerate}

\subsection{Child-Friendly Mental Model (Physical Metaphor)}
Imagine a team of 10 school children building an encyclopedia about their hometown:
\begin{itemize}
    \item \textbf{Without Redundancy Reduction}: All 10 children write an essay about the exact same local ice cream shop. The encyclopedia has 10 identical pages and is completely useless.
    \item \textbf{With Barlow Twins}:
    \begin{itemize}
        \item Child 1 must write about schools, Child 2 about parks, Child 3 about trains, Child 4 about rivers...
        \item If Child 1 tries to copy anything from Child 2, they get a penalty point (off-diagonal penalty).
        \item The resulting encyclopedia is rich, comprehensive, and contains zero duplicate information!
    \end{itemize}
\end{itemize}

\section{Core Mathematical Equations and Definitions}

\subsection{Cross-Correlation Matrix Construction}
{\boldmath
\begin{equation}
C_{ij} = \frac{1}{B} \sum_{b=1}^B \left( \frac{Z_{A, bi} - \mu_{A, i}}{\sigma_{A, i} + \epsilon} \right) \left( \frac{Z_{B, bj} - \mu_{B, j}}{\sigma_{B, j} + \epsilon} \right)
\end{equation}
}
\begin{itemize}
    \item \textbf{Formal Definition}: The empirical sample covariance matrix computed between the standardized feature coordinates of view $A$ and view $B$.
    \item \textbf{What It Means}: Measures how strongly coordinate $i$ in view $A$ co-varies with coordinate $j$ in view $B$ across the current mini-batch.
    \item \textbf{Why It Is Important}: Captures both cross-view consistency (diagonal) and cross-neuron redundancy (off-diagonal) in a single compact $D \times D$ operator.
\end{itemize}

\subsection{The On-Diagonal Invariance Term}
{\boldmath
\begin{equation}
\mathcal{L}_{\text{inv}} = \sum_{i=1}^D (1 - C_{ii})^2
\end{equation}
}
\begin{itemize}
    \item \textbf{Formal Definition}: The squared discrepancy penalizing diagonal correlation deviations from perfect unity.
    \item \textbf{What It Means}: Forces each individual feature dimension to remain perfectly invariant under data augmentations (cropping, color jitter, blurring).
    \item \textbf{Why It Is Important}: Ensures that the learned representations capture semantic content rather than spurious augmentation artifacts.
\end{itemize}

\subsection{The Off-Diagonal Redundancy Reduction Term}
{\boldmath
\begin{equation}
\mathcal{L}_{\text{red}} = \sum_{i=1}^D \sum_{j \ne i}^D C_{ij}^2
\end{equation}
}
\begin{itemize}
    \item \textbf{Formal Definition}: The sum of squared off-diagonal entries in the cross-correlation matrix.
    \item \textbf{What It Means}: Penalizes any linear correlation between different feature coordinates across views.
    \item \textbf{Why It Is Important}: Eliminates informational redundancy between neurons, forcing the network to utilize its full representational rank and preventing dimensional collapse.
\end{itemize}

\subsection{VICReg Variance Hinge Regularizer}
{\boldmath
\begin{equation}
v(\mathbf{Z}) = \frac{1}{D} \sum_{j=1}^D \max\left( 0, \, \gamma - \sqrt{\operatorname{Var}(Z_{\cdot, j}) + \epsilon} \right)
\end{equation}
}
\begin{itemize}
    \item \textbf{Formal Definition}: The hinge penalty enforcing that the empirical standard deviation of every feature coordinate across the batch exceeds margin $\gamma$ (typically $\gamma = 1$).
    \item \textbf{What It Means}: Directly forbids any neuron from becoming a static constant across different images.
    \item \textbf{Why It Is Important}: Provides a strict, explicit guarantee against complete representation collapse without requiring batch normalization.
\end{itemize}

\section{Rigorous Mathematical Analysis Checklist}

\subsection{Notation and Symbol Semantics}
\begin{itemize}
    \item $B \in \mathbb{N}$: Batch size ($B \ge 2$).
    \item $D \in \mathbb{N}$: Projector embedding dimension (typically $D \in [2048, 8192]$).
    \item $\mathbf{Z}_A, \mathbf{Z}_B \in \mathbb{R}^{B \times D}$: Output matrices of views $A$ and $B$.
    \item $\mathbf{C} \in \mathbb{R}^{D \times D}$: Empirical cross-correlation matrix.
    \item $\lambda \in \mathbb{R}^+$: Off-diagonal weight multiplier (typically $\lambda \approx 0.005$).
\end{itemize}

\subsection{Epistemic Status Table}
\begin{table}[h]
\centering
\small
\begin{tabular}{@{}llll@{}}
\toprule
\textbf{Quantity} & \textbf{Symbol} & \textbf{Epistemic Status} & \textbf{Operational Role} \\
\midrule
Cross-Correlation & $\mathbf{C}$ & Empirical Batch Statistic & Evaluates inter-view and inter-feature alignment \\
Invariance Penalty & $\mathcal{L}_{\text{inv}}$ & Exact Sum of Squares & Drives view invariance along diagonal \\
Redundancy Penalty & $\mathcal{L}_{\text{red}}$ & Exact Sum of Squares & Drives decorrelation across off-diagonal \\
Weight Multiplier & $\lambda$ & Tunable Hyperparameter & Balances invariance versus informational diversity \\
\bottomrule
\end{tabular}
\end{table}

\subsection{Computational Dependency Graph}
{\centering
$\{\mathbf{Z}_A, \mathbf{Z}_B\} \longrightarrow \text{Standardize along } B \longrightarrow \{\tilde{\mathbf{Z}}_A, \tilde{\mathbf{Z}}_B\} \longrightarrow \mathbf{C} = \frac{1}{B}\tilde{\mathbf{Z}}_A^T \tilde{\mathbf{Z}}_B \longrightarrow \{\mathcal{L}_{\text{inv}}, \mathcal{L}_{\text{red}}\} \longrightarrow \mathcal{L}_{\text{Barlow}}$
\par}

\subsection{Dimension and Coordinate Consistency}
The batch matrices have shape $B \times D$; the cross-correlation matrix $\mathbf{C}$ has shape $D \times D$. Matrix entries $C_{ij} \in [-1, 1]$ are dimensionless correlation coefficients.

\subsection{Asymptotic Regimes}
\begin{itemize}
    \item \textbf{Optimal Equilibrium}: $\mathbf{C} = \mathbf{I}_D \implies \mathcal{L}_{\text{Barlow}} = 0$; maximum mutual information with zero redundancy.
    \item \textbf{Complete Collapse}: All samples map to the same point $\implies C_{ij}$ is undefined (division by 0 standard deviation).
    \item \textbf{As $D \to \infty$}: Off-diagonal terms scale quadratically $\mathcal{O}(D^2)$; requires scaling $\lambda \propto 1/D$.
\end{itemize}

\subsection{Boundary Constraints and Invariants}
\begin{itemize}
    \item \textbf{Unit Diagonal Bound}: For any valid normalized vectors, $-1 \le C_{ii} \le 1$.
\end{itemize}

\subsection{Structural Isomorphisms}
Barlow Twins is structurally isomorphic to the classic Infomax principle (Linsker 1988) and Independent Component Analysis (ICA) applied across dual augmented views.

\section{Theoretical Theorems and Formal Proofs}

\begin{theorem}[Global Minimization of Barlow Twins Objective Implies Maximal Information Capacity]
Assume the representations $\mathbf{Z}_A, \mathbf{Z}_B$ follow a joint zero-mean multivariate Gaussian distribution $\mathcal{N}(\mathbf{0}, \mathbf{\Sigma})$. Then the global minimum $\mathcal{L}_{\text{Barlow}} = 0$ is achieved if and only if $\mathbf{C} = \mathbf{I}_D$, which maximizes the mutual information $I(\mathbf{Z}_A; \mathbf{Z}_B)$ subject to an isotropic orthogonal representation budget.
\end{theorem}

\begin{proof}
Consider the objective function $\mathcal{L}_{\text{Barlow}} = \sum_{i=1}^D (1 - C_{ii})^2 + \lambda \sum_{i=1}^D \sum_{j \ne i}^D C_{ij}^2$.
Since $\lambda > 0$ and all terms are squared real numbers:
{\boldmath
\begin{equation}
\mathcal{L}_{\text{Barlow}} \ge 0
\end{equation}
}
Equality $\mathcal{L}_{\text{Barlow}} = 0$ holds if and only if each term in the summation vanishes individually:
{\boldmath
\begin{align}
(1 - C_{ii})^2 = 0 &\iff C_{ii} = 1 \quad \forall i \in \{1, \dots, D\} \\
C_{ij}^2 = 0 &\iff C_{ij} = 0 \quad \forall i \ne j
\end{align}
}
This is identically the condition that the cross-correlation matrix $\mathbf{C}$ equals the identity matrix $\mathbf{I}_D$:
{\boldmath
\begin{equation}
\mathbf{C} = \mathbf{I}_D
\end{equation}
}
For jointly Gaussian vectors $(\mathbf{z}_A, \mathbf{z}_B)$ with covariance $\operatorname{Cov}(\mathbf{z}_A) = \operatorname{Cov}(\mathbf{z}_B) = \mathbf{I}_D$ and cross-covariance $\mathbf{C}$, the mutual information is:
{\boldmath
\begin{equation}
I(\mathbf{z}_A; \mathbf{z}_B) = -\frac{1}{2} \log \det\left( \mathbf{I}_D - \mathbf{C}^T \mathbf{C} \right)
\end{equation}
}
As $C_{ii} \to 1$ and $C_{ij} \to 0$ ($i \ne j$), each coordinate pair $i$ achieves maximal mutual information $I(z_{A, i}; z_{B, i}) \to \infty$ independently, while mutual information across disparate coordinates $I(z_{A, i}; z_{B, j}) = 0$ for $i \ne j$.
Thus, achieving $\mathcal{L}_{\text{Barlow}} = 0$ maximally retains task-invariant information while eliminating redundant inter-feature crosstalk.
\end{proof}

\begin{theorem}[Robustness to Mini-Batch Size Scaling]
Unlike contrastive estimators whose bound on mutual information is constrained by $\log(B)$, the gradient magnitude of the Barlow Twins loss per feature parameter does not diminish with small batch sizes $B$, provided $B > 1$.
\end{theorem}

\begin{proof}
In contrastive InfoNCE loss, the denominator sums over $B$ samples: $\mathcal{L}_{\text{NCE}} = -\log \frac{\exp(\mathbf{z}_i \cdot \mathbf{z}_i^+ / \tau)}{\sum_{b=1}^B \exp(\mathbf{z}_i \cdot \mathbf{z}_b / \tau)}$. The mutual information estimate is mathematically bounded by $I \le \log(B)$. If $B$ is small, the gradient cannot distinguish fine negative structures.

In Barlow Twins, the loss is computed over the $D \times D$ correlation matrix $\mathbf{C}$.
The gradient of $C_{ij}$ with respect to activation $\tilde{Z}_{A, bi}$ is:
{\boldmath
\begin{equation}
\frac{\partial C_{ij}}{\partial \tilde{Z}_{A, bi}} = \frac{1}{B} \tilde{Z}_{B, bj}
\end{equation}
}
By the chain rule:
{\boldmath
\begin{equation}
\frac{\partial \mathcal{L}}{\partial \tilde{Z}_{A, bi}} = -\frac{2}{B}(1 - C_{ii})\tilde{Z}_{B, bi} + \frac{2\lambda}{B} \sum_{j \ne i} C_{ij} \tilde{Z}_{B, bj}
\end{equation}
}
Summing the gradient norms across the batch:
{\boldmath
\begin{equation}
\sum_{b=1}^B \left( \frac{\partial \mathcal{L}}{\partial \tilde{Z}_{A, bi}} \right)^2 \propto \mathcal{O}\left( (1 - C_{ii})^2 + \lambda^2 \sum_{j \ne i} C_{ij}^2 \right)
\end{equation}
}
Notice that the total gradient energy is independent of batch size $B$ and depends entirely on the distance from the identity matrix $\mathbf{I}_D$. Thus, representation decorrelation proceeds effectively even for modest batch sizes.
\end{proof}

\section{Foundational Architectural Questions and Epistemic Meta-Questions}

\subsection{Architectural Question 1: Dimensionality of the Projector Network}
\textbf{Question:} Why does Barlow Twins require very large projector dimensions (e.g. $D = 8192$) compared to standard contrastive learning ($D = 128$)?\\
\textbf{Answer:} Barlow Twins relies on the cross-correlation matrix $\mathbf{C} \in \mathbb{R}^{D \times D}$. Larger dimensions provide more orthogonal coordinate axes to disentangle independent semantic factors. Increasing $D$ from 256 to 8192 increases linear probing performance monotonically because it forces the network to discover hundreds of non-overlapping visual concepts.

\textbf{Meta-Question:} Doesn't an $8192 \times 8192$ correlation matrix create an excessive computational bottleneck?\\
\textbf{Answer:} No. The matrix multiplication $\tilde{\mathbf{Z}}_A^T \tilde{\mathbf{Z}}_B$ is computed only once per mini-batch in the low-dimensional projection head, requiring $\mathcal{O}(B \cdot D^2)$ operations, which is negligible compared to the billions of FLOPs spent inside the vision backbone.

\subsection{Architectural Question 2: Barlow Twins vs VICReg}
\textbf{Question:} How does VICReg improve upon Barlow Twins for non-image or tabular modalities?\\
\textbf{Answer:} Barlow Twins requires batch normalization along feature dimensions to standardize variance. In small-batch or non-stationary domains (e.g. graph or tabular data), batch normalization can be unstable. VICReg replaces batch normalization with an explicit variance hinge loss $v(\mathbf{Z}) = \max(0, 1 - \operatorname{std}(Z))$, allowing completely batch-norm-free self-supervised pre-training.

\textbf{Meta-Question:} Does the variance hinge loss in VICReg introduce non-smooth gradients?\\
\textbf{Answer:} The hinge function has a subgradient at the threshold $\gamma$, but in practice feature variances start well below $\gamma = 1$, generating smooth linear gradients until stability is reached.

\section{Algorithmic Implementation Specification}

\begin{algorithm}[H]
\caption{Barlow Twins Batch Standardization and Redundancy-Reduction Loss}
\begin{algorithmic}[1]
\Require Batch representation matrices $\mathbf{Z}_A, \mathbf{Z}_B \in \mathbb{R}^{B \times D}$, parameter $\lambda \ge 0$
\Ensure Total loss $\mathcal{L}_{\text{Barlow}}$, on-diagonal loss $\mathcal{L}_{\text{inv}}$, off-diagonal loss $\mathcal{L}_{\text{red}}$, mean diagonal correlation $\bar{C}_{\text{diag}}$
\State $\epsilon \gets 10^{-8}$
\State Initialize normalized matrices $\tilde{\mathbf{Z}}_A, \tilde{\mathbf{Z}}_B$ of shape $B \times D$
\For{$j = 0$ \textbf{to} $D - 1$}
    \State $\mu_A \gets \frac{1}{B}\sum_{b=0}^{B-1} \mathbf{Z}_A[b][j]$, \quad $\sigma_A \gets \sqrt{\frac{1}{B}\sum_{b=0}^{B-1} (\mathbf{Z}_A[b][j] - \mu_A)^2} + \epsilon$
    \State $\mu_B \gets \frac{1}{B}\sum_{b=0}^{B-1} \mathbf{Z}_B[b][j]$, \quad $\sigma_B \gets \sqrt{\frac{1}{B}\sum_{b=0}^{B-1} (\mathbf{Z}_B[b][j] - \mu_B)^2} + \epsilon$
    \For{$b = 0$ \textbf{to} $B - 1$}
        \State $\tilde{\mathbf{Z}}_A[b][j] \gets (\mathbf{Z}_A[b][j] - \mu_A) / \sigma_A$
        \State $\tilde{\mathbf{Z}}_B[b][j] \gets (\mathbf{Z}_B[b][j] - \mu_B) / \sigma_B$
    \EndFor
\EndFor
\State $\mathcal{L}_{\text{inv}} \gets 0.0$, \quad $\mathcal{L}_{\text{red}} \gets 0.0$, \quad $diag\_sum \gets 0.0$
\For{$i = 0$ \textbf{to} $D - 1$}
    \For{$j = 0$ \textbf{to} $D - 1$}
        \State $C_{ij} \gets \frac{1}{B}\sum_{b=0}^{B-1} \tilde{\mathbf{Z}}_A[b][i] \cdot \tilde{\mathbf{Z}}_B[b][j]$
        \If{$i == j$}
            \State $\mathcal{L}_{\text{inv}} \gets \mathcal{L}_{\text{inv}} + (1.0 - C_{ij})^2$
            \State $diag\_sum \gets diag\_sum + C_{ij}$
        \Else
            \State $\mathcal{L}_{\text{red}} \gets \mathcal{L}_{\text{red}} + C_{ij}^2$
        \EndIf
    \EndFor
\EndFor
\State $\mathcal{L}_{\text{Barlow}} \gets \mathcal{L}_{\text{inv}} + \lambda \cdot \mathcal{L}_{\text{red}}$
\State \Return $(\mathcal{L}_{\text{Barlow}}, \mathcal{L}_{\text{inv}}, \mathcal{L}_{\text{red}}, diag\_sum / D)$
\end{algorithmic}
\end{algorithm}

\section{Big-$\Theta$ Complexity Analysis}

\begin{itemize}
    \item \textbf{Batch Normalization Time Complexity}: $\Theta(B \cdot D)$ scalar operations across the mini-batch.
    \item \textbf{Cross-Correlation Matrix Multiplications Complexity}: $\Theta(B \cdot D^2)$ operations to compute $\mathbf{C} = \frac{1}{B}\tilde{\mathbf{Z}}_A^T \tilde{\mathbf{Z}}_B$.
    \item \textbf{Space Complexity}: $\Theta(B \cdot D + D^2)$ memory for standardized representations and correlation buffers.
    \item \textbf{Work \& Depth}: Work is $\mathcal{W} = \Theta(B \cdot D^2)$; parallel depth across matrix contractions is $\mathcal{D} = \Theta(\log B + \log D)$.
\end{itemize}

\section{Single Canonical Repository Reference}

The complete deterministic scalar implementation, verification suite, and unit test benchmarks are published at:
\begin{center}
\url{https://github.com/Chief-Strategist-J/llm-obs-infra}
\end{center}

\end{document}
